% GENERATED FILE. Edit canonical lessons, then run npm run generate:books.
% canonical-source-hash: fnv1a64:e009f00e4ca25c82
% canonical-lessons: JA-C62-nimotsu, JA-C62-fune, JA-C62-oto, JA-C62-iro, JA-C62-katachi

\chapter{Luggage, Ship, Sound, Colour, Shape}
\label{ch:ja-luggage-ship-sound-colour-shape}
\hlchaptermodality{62}

\begin{chapteropening}
\textbf{By the end of this chapter:} \emph{I can say luggage, a ship, a sound, a colour and a shape.}

Complete the last lesson of chapter 62: i can say luggage, a ship, a sound, a colour and a shape.
\end{chapteropening}

\section[nimotsu]{\ja{にもつ} (nimotsu) — luggage}

\label{lesson:JA-C62-nimotsu}



\begin{quote}
Before the new one: say the Japanese for a string, then the Japanese for a rope.
\end{quote}

\subsection*{You'll want to know: \ja{にもつ}}
\textbf{\ja{にもつ}} — \emph{nimotsu} — \textquotedblleft{}luggage\textquotedblright{}.

\begin{grammarlens}[title={What you've built}]
One more word for this chapter.
\end{grammarlens}

\subsection*{Your turn}
\begin{itemize}
\raggedright
  \item \emph{Say it:} \emph{nimotsu}
  \item \emph{Say it:} \emph{nimotsu}, once more
  \item \emph{Say it:} say \emph{nimotsu}
  \item \emph{Recall:} say the Japanese for a string, then the Japanese for a rope, then say \emph{nimotsu} again
\end{itemize}

\begin{tcolorbox}[breakable,colback=teal!4,colframe=teal!35!black,title={Before you move on}]
What does \ja{にもつ} mean? (\textquotedblleft{}Luggage\textquotedblright{}.) Say it once more.
\end{tcolorbox}
\section[fune]{\ja{ふね} (fune) — a ship}

\label{lesson:JA-C62-fune}



\begin{quote}
Before the new one: say the Japanese for a rope, then the Japanese for luggage.
\end{quote}

\subsection*{You'll want to know: \ja{ふね}}
\textbf{\ja{ふね}} — \emph{fune} — \textquotedblleft{}a ship\textquotedblright{}.

\begin{grammarlens}[title={What you've built}]
One more word for this chapter.
\end{grammarlens}

\subsection*{Your turn}
\begin{itemize}
\raggedright
  \item \emph{Say it:} \emph{fune}
  \item \emph{Say it:} \emph{fune}, once more
  \item \emph{Say it:} say \emph{fune}
  \item \emph{Recall:} say the Japanese for a rope, then the Japanese for luggage, then say \emph{fune} again
\end{itemize}

\begin{tcolorbox}[breakable,colback=teal!4,colframe=teal!35!black,title={Before you move on}]
What does \ja{ふね} mean? (\textquotedblleft{}A ship\textquotedblright{}.) Say it once more.
\end{tcolorbox}
\section[oto]{\ja{おと} (oto) — a sound}

\label{lesson:JA-C62-oto}



\begin{quote}
Before the new one: say the Japanese for luggage, then the Japanese for a ship.
\end{quote}

\subsection*{You'll want to know: \ja{おと}}
\textbf{\ja{おと}} — \emph{oto} — \textquotedblleft{}a sound\textquotedblright{}.

\begin{grammarlens}[title={What you've built}]
One more word for this chapter.
\end{grammarlens}

\subsection*{Your turn}
\begin{itemize}
\raggedright
  \item \emph{Say it:} \emph{oto}
  \item \emph{Say it:} \emph{oto}, once more
  \item \emph{Say it:} say \emph{oto}
  \item \emph{Recall:} say the Japanese for luggage, then the Japanese for a ship, then say \emph{oto} again
\end{itemize}

\begin{tcolorbox}[breakable,colback=teal!4,colframe=teal!35!black,title={Before you move on}]
What does \ja{おと} mean? (\textquotedblleft{}A sound\textquotedblright{}.) Say it once more.
\end{tcolorbox}
\section[iro]{\ja{いろ} (iro) — a colour}

\label{lesson:JA-C62-iro}



\begin{quote}
Before the new one: say the Japanese for a ship, then the Japanese for a sound.
\end{quote}

\subsection*{You'll want to know: \ja{いろ}}
\textbf{\ja{いろ}} — \emph{iro} — \textquotedblleft{}a colour\textquotedblright{}.

\begin{grammarlens}[title={What you've built}]
One more word for this chapter.
\end{grammarlens}

\subsection*{Your turn}
\begin{itemize}
\raggedright
  \item \emph{Say it:} \emph{iro}
  \item \emph{Say it:} \emph{iro}, once more
  \item \emph{Say it:} say \emph{iro}
  \item \emph{Recall:} say the Japanese for a ship, then the Japanese for a sound, then say \emph{iro} again
\end{itemize}

\begin{tcolorbox}[breakable,colback=teal!4,colframe=teal!35!black,title={Before you move on}]
What does \ja{いろ} mean? (\textquotedblleft{}A colour\textquotedblright{}.) Say it once more.
\end{tcolorbox}
\section[katachi]{\ja{かたち} (katachi) — a shape}

\label{lesson:JA-C62-katachi}



\begin{quote}
Before the new one: say the Japanese for a sound, then the Japanese for a colour.
\end{quote}

\subsection*{You'll want to know: \ja{かたち}}
\textbf{\ja{かたち}} — \emph{katachi} — \textquotedblleft{}a shape\textquotedblright{}.

\begin{grammarlens}[title={What you've built}]
One more word for this chapter.
\end{grammarlens}

\subsection*{Your turn}
\begin{itemize}
\raggedright
  \item \emph{Say it:} \emph{katachi}
  \item \emph{Say it:} \emph{katachi}, once more
  \item \emph{Say it:} say \emph{katachi}
  \item \emph{Recall:} say the Japanese for a sound, then the Japanese for a colour, then say \emph{katachi} again
\end{itemize}

\begin{tcolorbox}[breakable,colback=teal!4,colframe=teal!35!black,title={Before you move on}]
What does \ja{かたち} mean? (\textquotedblleft{}A shape\textquotedblright{}.) Say it once more.
\end{tcolorbox}
